% This must be in the first 5 lines to tell arXiv to use pdfLaTeX, which is strongly recommended.
\pdfoutput=1
% In particular, the hyperref package requires pdfLaTeX in order to break URLs across lines.

\documentclass[11pt]{article}

% Remove the "review" option to generate the final version.
\usepackage{ACL2023}
%\usepackage{ACL2023}

% Standard package includes
\usepackage{times}
\usepackage{latexsym,comment}

% For proper rendering and hyphenation of words containing Latin characters (including in bib files)
\usepackage[T1]{fontenc}
% For Vietnamese characters
% \usepackage[T5]{fontenc}
% See https://www.latex-project.org/help/documentation/encguide.pdf for other character sets

% This assumes your files are encoded as UTF8
\usepackage[utf8]{inputenc}

% This is not strictly necessary, and may be commented out.
% However, it will improve the layout of the manuscript,
% and will typically save some space.
\usepackage{microtype}

% This is also not strictly necessary, and may be commented out.
% However, it will improve the aesthetics of text in
% the typewriter font.
\usepackage{inconsolata}


\usepackage{natbib}
\usepackage{graphicx}
\usepackage{appendix}
\usepackage{microtype}
\usepackage{bm}
\usepackage{multirow}
\usepackage{enumitem}
\usepackage{soul} % for the command \hl
\usepackage{times}
\usepackage{latexsym}
\usepackage{graphicx}
\usepackage{booktabs}
\usepackage{tcolorbox}
\usepackage{paralist}
\usepackage{enumitem}
%\setlist{topsep=0pt, leftmargin=*}

\usepackage{mathtools}
\usepackage{amsfonts,amssymb,amsthm}
\usepackage[ruled,noend]{algorithm2e}


\usepackage{cleveref}
\crefformat{section}{\S#2#1#3}
\crefformat{subsection}{\S#2#1#3}
\crefformat{subsubsection}{\S#2#1#3}
\crefrangeformat{section}{\S\S#3#1#4 to~#5#2#6}
\crefmultiformat{section}{\S\S#2#1#3}{ and~#2#1#3}{, #2#1#3}{ and~#2#1#3}


%\usepackage{refstyle}
\Crefformat{figure}{#2Figure~#1#3}
\Crefmultiformat{figure}{Figures~#2#1#3}{ and~#2#1#3}{, #2#1#3}{ and~#2#1#3}
\Crefformat{table}{#2Table~#1#3}
\Crefmultiformat{table}{Tables.~#2#1#3}{ and~#2#1#3}{, #2#1#3}{ and~#2#1#3}
\Crefformat{appendix}{#2Appx.~\S#1#3}
\crefformat{algorithm}{Alg.~#2#1#3}
\Crefformat{equation}{#2Eq.~#1#3}
\newcommand{\stitle}[1]{\vspace{0.3em}\noindent{\bf #1}}



\pagestyle{plain} % page number

\newcommand{\model}{\textsc{MatCha}}

% \usepackage{dblfloatfix}
\usepackage{color, colortbl}
\usepackage{xcolor}
\definecolor{asparagus}{rgb}{0.53, 0.66, 0.42}
\definecolor{applegreen}{rgb}{0.55, 0.71, 0.0}

\newcommand{\todo}[1]{\textcolor{red}{{\bf TODO: #1}}}

\newif\ifcomment\commenttrue
\ifcomment
\newcommand{\pinaforecomment}[3]{\colorbox{#1}{\parbox{.8\linewidth}{#2: #3}}}
\else
\newcommand{\pinaforecomment}[3]{}
\fi

\newcommand{\flcomment}[1]{\pinaforecomment{red}{FL}{#1}}
\newcommand{\juliancomment}[1]{\pinaforecomment{green}{JE}{#1}}
% orange, green, etc

\title{\model{} : Enhancing Visual Language Pretraining \\with Math Reasoning and Chart Derendering}
\author{Fangyu Liu$^{\spadesuit\clubsuit}$\thanks{\ \ Work done during Google internship.} \ \ \ Francesco Piccinno$^\clubsuit$ \ \ \ Syrine Krichene$^\clubsuit$ \ \ \ Chenxi Pang$^\clubsuit$ \ \ \ Kenton Lee$^\clubsuit$ \\ \textbf{Mandar Joshi$^\clubsuit$ \ \ \ Yasemin Altun$^\clubsuit$ \ \ \ Nigel Collier$^\spadesuit$ \ \ \ Julian Martin Eisenschlos$^\clubsuit$} \\
$^\clubsuit$Google DeepMind \ \ \ \ $^\spadesuit$University of Cambridge}

\begin{document}

\maketitle


\begin{abstract}
Visual language data such as plots, charts, and infographics are ubiquitous in the human world. However, state-of-the-art vision-language models do not perform well on these data.
We propose \model{} (\textbf{Mat}h reasoning and \textbf{Cha}rt derendering pretraining) to enhance visual language models' capabilities in jointly modeling charts/plots and language data. Specifically we propose several pretraining tasks that cover plot deconstruction and numerical reasoning which are the key capabilities in visual language modeling. 
% and plot deconstruction capabilities. These two capabilities enable a model with the key capabilities of (1) extracting key information and (2) reasoning on the extracted information. 
We perform the \model{} pretraining starting from Pix2Struct, a recently proposed image-to-text visual language model.
% and continue pretraining with our proposed objectives. 
On standard benchmarks such as PlotQA and ChartQA, the \model{} model outperforms state-of-the-art methods by as much as nearly 20\%.
We also examine how well the \model{} pretraining transfers to domains such as screenshots, textbook diagrams, and document figures and observe overall improvement, verifying the usefulness of \model{} pretraining on broader visual language tasks.\footnote{Code and models: \href{https://github.com/google-research/google-research/tree/master/deplot}{github.com/google-research/google-research/tree/master/deplot}}\footnote{For questions paper please contact \texttt{fl399@cam.ac.uk} and \texttt{eisenjulian@google.com}.}
%We observe improvement over the base Pix2Struct checkpoint by 1.2\% on average, 

\end{abstract}

\section{Introduction}\label{sec:intro}

Visual language is the system that uses tightly integrated textual and visual elements to convey meaning \citep{horn1998visual}.
It is ubiquitous in the human world with typical examples being charts, plots and diagrams existing in places such as textbooks, scientific papers web pages and many more. % infographics etc.
Visual language is also highly complex -- besides texts, its structural units can include line, shape, color, orientation, scale, angle, space, etc. One needs to recognize patterns from these structural units, and perform spatial grouping and/or alignment to extract information for reasoning.

Whilst being prevalent and important, there is little research on visual language understanding from the machine learning community.
%has been largely overlooked by the machine learning community.
Vision-language models pretrained on natural images or image-text pairs crawled from the web perform badly on visual language tasks such as ChartQA \citep{masry-etal-2022-chartqa} and PlotQA \citep{methani2020plotqa} due to the high complexity of jointly modeling language and symbols (more evidence in experiments). Pix2Struct \citep{lee2022pix2struct} is a recently proposed pretraining strategy for visually-situated language that significantly outperforms standard vision-language models, and also a wide range of OCR-based pipeline approaches. Pix2Struct designs a novel masked webpage screenshot parsing task and also a variable-resolution input representation for pretraining an image-to-text encode-decoder Transformer \citep{vaswani2017attention}. In this work, we use Pix2Struct as the base model and further pretrain it with chart derendering and math reasoning tasks.

\begin{figure*}[!ht]
\centering
\includegraphics[width=\linewidth]{figs/matcha_front_page_fig.pdf}
\caption{\model{} defines two types of pretraining tasks: (1) chart derendering (light blue boxes) and (2) mathematical reasoning (light red boxes). In chart derendering, given a chart, the model needs to decode its underlying rendering code or data table. In math reasoning, given a math question rendered as an image, the model needs to decode its answer. Chart derendering teaches the models layout understanding (including number extraction and their organizations) and math reasoning teaches the models numerical reasoning capabilities. }
\label{fig:front_page}
\end{figure*}

We argue that visual language understanding needs two key ingredients: (1) layout understanding (including number extraction and their organizations) and (2) mathematical reasoning. (1) is required to discover the underlying patterns of the image and organize the elements in the image in a logical form. (2) is needed to operate on the elements extracted from (1) and derive meaningful information demanded by a task or query. Based on these observations, we propose two complementary pretraining tasks for enhancing visual language understanding: \textbf{chart derendering} and \textbf{math reasoning}. In chart derendering, given a plot/chart, the image-to-text model is required to generate its underlying data table or the code used to render it. The second task is math reasoning pretraining. We pick two numerical reasoning dataset MATH \citep{saxton2019analysing} and DROP \citep{dua-etal-2019-drop}, render the input into images and the image-to-text model needs to decode the answers.


We use a suite of visual language tasks to test the effectiveness of our method. Most importantly, we test on ChartQA and PlotQA which are QA datasets about plots and charts. On both datasets, \model{} surpasses even the SOTA model assuming access to charts' underlying data tables and can beat the prior SOTA without gold data table by as much as 20\%. We also test \model{} on chart-to-text summarization tasks and observe clear improvements over Pix2Struct and achieves SOTA on Chart-to-Text \citep{kantharaj-etal-2022-chart} Pew split. Last but not least, to examine if the \model{} pretraining generalizes to datasets beyond the standard plots and charts domain, we also test \model{} on four additional domains where Pix2Struct was evaluated on: documents, illustrations, user
interfaces, and natural images (including datasets, such as textbook QA, Widget Captioning, etc.).
We demonstrate consistent improvement on most additional datasets compared with the base model Pix2Struct. 

To summarize, our contributions are: (1) proposing a set of effective pretraining tasks for visual language learning (2) demonstrating consistent improvements across all evaluated tasks and SOTA results on ChartQA, PlotQA, and Chart-to-Text summarization (Statista set) without accessing the gold data tables; (3) verify that \model{} pretraining transfers to visual language benchmarks beyond the chart \& plot domains and achieve SOTA across a wide range of datasets beyond the chart domain such as textbook VQA and Widget Captioning; (4) comprehensive ablation and analyses to understand the effect of each pretraining component and its impact to downstream performance.


\section{Related Work}\label{sec:rw}

\paragraph{Vision-language research and a lack of attention on visual language.}
%Better define the concept. Put it into the right context.
Research on vision-and-language has predominately been focusing on natural images. Visually-grounded reasoning datasets such as NLVR2 \citep{suhr-etal-2019-corpus} and MaRVL \citep{liu-etal-2021-visually} are mostly in the natural image domain. Synthesized datasets such as SHAPES \citep{andreas2016neural}, NLVR \citep{suhr-etal-2017-corpus}, and CLEVR \citep{johnson2017clevr} can be seen as in the visual language domain. However, their visual language systems are significantly simpler than those in the real world such as plots and charts. As a result, information extraction from these synthesized datasets is straightforward. Besides, the queries in the synthesized datasets are relatively naive and do not require complex reasoning (e.g., questions can usually be on spatial relations or counting objects). Consequently, current vision-language models can handle the above mentioned synthesized visual reasoning datasets quite well. However, they do not perform well on real-world visual language datasets where both the information extraction and reasoning becomes much more complex (we will show this in \Cref{sec:exp}).

\paragraph{OCR-based \& end-to-end methods for visually-situated language.\footnote{There is a nuanced difference between \emph{visual} language and \emph{visually-situated} language. Most models discussed here are specifically designed for images with significant amount of texts (e.g., documents) and thus they are models primarily for visually-situated language. Visual language data significantly overlaps with visually-situated language data. However, visual language also covers the scenarios where no/few texts are explicitly used but visual objects/patterns are most responsible for presenting information (e.g., certain types of plots).}}
LayoutLM \citep{xu2020layoutlm,huang2022layoutlmv3} leverages a patch-OCR alignment loss to inject an external OCR systems' knowledge into the Transformer model.
PresSTU \citep{kil2022prestu} and PaLI \citep{chen2022pali} also design OCR-aware pretraining objectives where the model needs to predict texts obtained from off-the-shelf OCR systems. 
ChartBERT \citep{akhtar-etal-2023-reading} relies on OCR text and positions to train a transformer encoder.
While OCR systems can be helpful for accurately extracting texts, running them is not cheap. Also, OCR systems do not cover visual language systems that do not explicitly use text. As examples, plots and charts do not always have numbers written explicitly.
In our concurrent work \textsc{DePlot} \citep{liu2022deplot}, we explore combining a chart-to-text translation module (without OCR) with large language models.

%\paragraph{End-to-end methods for visual language.}
Donut \citep{kim2022ocr}, Dessurt \citep{Davis2022EndtoendDR}, and Pix2Struct \citep{lee2022pix2struct} are end-to-end pretrained models for visual language where Donut and Dessurt focus on document understanding and Pix2Struct aim to provide a generic pretrained checkpoint for all visual language tasks. \model's architecture is identical to Pix2Struct -- we continually pretrain a Pix2Struct checkpoint with new objectives.

\paragraph{Learning to reason by designing novel pretraining tasks.} \model{} is related to the literature of designing better pretraining objectives to help the language models (LMs) to reason better since the skill is hard to require through naive language modeling objectives only (e.g, masked language modeling and autoregressive language modeling on raw text).
\citet{geva-etal-2020-injecting, eisenschlos-etal-2020-understanding} generate additional pretraining data focused on (numerical) reasoning through human-written templates. 
\citet{pi-etal-2022-reasoning} synthesize data and programs, and then use program executors to simulate answers. LMs are pretrained to predict the answers given data and programs.
%TAPEX \citep{liu2022tapex} 
\citet{wu2022insights} explore a wide range of synthetic pretraining tasks and found that even just injecting knowledge as simple as induction and deduction rules could teach LMs to reason. We teach an image-to-text model to reason through mapping charts to data and code, and also directly learning textual math reasoning datasets.


\section{Method}\label{sec:method}

We argue that layout understanding and basic math operation capabilities are the key elements for performing visual language understanding/reasoning. We inject such capabilities to the model by proposing two pretraining tasks: \textbf{chart derendering} (\Cref{sec:plot_derendering}) and \textbf{math reasoning} (\Cref{sec:math_reasoning}) which we introduce in detail in the following sections.

\subsection{Chart Derendering}\label{sec:plot_derendering}
Plots and charts are usually generated by an underlying data table and a piece of code. Code decides the overall layout of the figure (e.g., type, direction, color/shape scheme of the chart) and the underlying data table decides the actual numbers and the groupings of them. Both the data and code are sent to a compiler/rendering engine to create the final image. 
To understand a chart one needs to discover the visual patterns in the image, effectively parse and group them to extract the key information. Reversing the plot rendering process demands all such capabilities and can thus serve as a perfect pretraining task.

In practice, it is challenging to simultaneously obtain charts, their underlying data tables, and their rendering code. To collect sufficient pretraining data, we independently accumulate (chart, code) and (chart, table) pairs. For (chart, code), we crawl all GitHub IPython notebooks with appropriate licenses and extract blocks with figures. A figure and the code block right before it are saved as a (chart, code) pair.\footnote{Note that the code snippet can be noisy since earlier blocks could also be relevant for generating the figure and also the snippet may contain bits of code that is irrelevant to generating the figure. Also note that the data table is frequently missing and usually not hardcoded in the notebook. As a result, we collect (chart, table) pairs separately.} For (chart, table) pairs, we explored two sources. First is to manually write code for converting web-crawled Wikipedia tables from ~\citet{herzig-etal-2020-tapas} to charts.
%around 200 lines of Python code, 
We randomly combine several plotting options. The key random variables include: using either \texttt{matplotlib} or \texttt{seaborn} as the plotting package; using either bar, line, or pie chart; styles and colors of the charts; whether to show numbers explicitly on the graph; font and size of the texts. Besides our own synthetic data, we also add chart-table pairs generated by %\citet{masry-etal-2022-chartqa} and
\citet{methani2020plotqa} (from PlotQA) to diversify the pretraining corpus. The second source is web-crawled chart-table pairs. Websites such as Statista provides both. We directly use the chart-table pairs crawled by \citet{masry-etal-2022-chartqa} (from ChartQA), containing around 20k pairs in total from four websites: Statista, Pew, Our World in Data, and OECD.\footnote{See \Cref{sec:dataset_details} for links.}

Note that to avoid leaking test information for the PlotQA and ChartQA tasks which use the same chart data as pretraining, we only use the chart-table pairs in the training sets for pretraining and test tables/charts are strictly excluded. In ablation study (\Cref{sec:ablations}), we will show that chart-table from both sources are useful and having a diverse set of chart-table pairs is always better. However, using only our synthetic data brings very significant improvement already, suggesting that the concept of chart derendering can be easily transferred to charts of other domains (including real-world charts).

%TODO: show some pretraining examples for the tasks


\subsection{Math Reasoning}\label{sec:math_reasoning}
Reasoning over visual language requires (1) effective recognition and grouping of the visual elements and also (2) applying mathematical operations (such as sorting, min/max, averaging, etc.) on top of them. Plot derendering addresses (1) but (2) is still lacking in the current pretraining framework. As a result, we propose to explicitly inject numerical reasoning knowledge to the image-to-text model by learning from textual math datasets.

We use two existing textual math reasoning datasets, MATH \citep{saxton2019analysing} and DROP \citep{dua-etal-2019-drop} for pretraining. MATH is synthetically created, containing two million training examples per module (type) of questions (see \Cref{sec:dataset_details} for a comprehensive listing of modules included in \model{} pretraining).
DROP is a reading-comprehension-style QA dataset where the input is a paragraph context and a question. DROP has 96k question and answer pairs over 6.7K paragraphs.\footnote{Note that for all datasets used for pretraining, we always use only the training set if there exists a split.} To solve questions in DROP, the model needs to read the paragraph, extract relevant numbers and perform numerical computation to predict the answer. We found both datasets to be complementarily helpful. MATH contains large amounts of questions and is categorized which helps us identify math operations needed to explicitly inject to the model. DROP's reading-comprehension format resembles the typical QA format where models need to simultaneously perform information extraction and reasoning. In practice, we render inputs of both datasets into images (concatenating the context and question for DROP). The image-to-text model is trained to decode the answer given the redered image. Examples of MATH and DROP can be found in \Cref{fig:front_page} (in light red).

Besides the two newly proposed pretraining strategies, to prevent catastrophic forgetting, we also keep applying the screenshot parsing pretraining from Pix2Struct \citep{lee2022pix2struct}. Specifically, given screenshot of a website (where parts of the website is masked), the image-to-text transformer needs to predict the underlying simplified HTML code that could render the original unmasked website screenshot. The final pretraining task is a mixture of all aforementioned tasks. We discuss the mixture weights in \Cref{sec:exp_setups}.

\input{04_experiment}
\section{Conclusion}

We have proposed a pretraining method \model{} for visual language tasks. \model{} injects chart understanding and reasoning knowledge to an image-to-text transformer model by learning to (1) predict the underlying data tables and code given chart images and (2) decode the answers of math questions (rendered in the form of images). \model{} establishes new SOTA on 5 out of 6 setups across three chart domain benchmarks covering both QA and summarization tasks. On visual language tasks beyond the chart domain (e.g., textbook QA and DocVQA), \model{} improves upon Pix2Struct, indicating that the learned knowledge in \model{} pretraining can be transferred outside of the pretraining domain. We conduct comprehensive ablation studies to identify the actual impact of each pretraining component and task and find that chart derendering is essential for extractive questions while math pretraining is important for queries that requires complex reasoning.

\section*{Limitations}\label{sec:limits}
Though we have injected math reasoning skills to \model, error analysis shows that there is still room for improvement on queries requiring complex reasoning. Besides, it remains debatable whether doing math calculation in weight space in a purely end-to-end manner is the most promising path forward.\footnote{See recent works that combine LLMs with calculators \citep{wei2022chain}  or compilers/program executors \citep{cheng2022binding,chen2022program,gao2022pal}.}

Besides math reasoning, \Cref{fig:finegrained} shows that plot attributes is an area where \model{} underperforms PaLI. We conjecture that it is due to \model's lack of massive scale grounded image-text pretraining with rich semantics (which PaLI has using web-scale image-text pairs). While chart-to-code pretraining provides certain level of plot attribute grounding, such plot features are mostly using default options in plotting packages but not explicitly written out in code.
%We plan to improve \model{} on this front in future work.


In terms of experimental setup, the reported number is result of a single run. Pretraining is extremely costly especially when there exists more than twenty ablation setups and downstream evaluation tasks. We have collected pretraining and evaluation data points from multiple aspects on various scenarios to verify the robustness of \model. However, we do acknowledge that the paper can benefit from reporting multiple runs given sufficient compute.


Last but not least, it is also worth noting that visual language is an umbrella term. There are other visual language systems beyond the ones discussed in this paper. As an example, comics/manga have their distinct visual lexicon or even grammars \citep{cohn2013visual}.

\section*{Ethics Statement}\label{sec:ethics}
To the best of our knowledge, \model{} has not been trained on sensitive private information and should be of low risk to generate harmful contents. All pretraining and finetuning data are either synthetically created using rules or publicly available data on the web with appropriate permissive licenses.

% \section*{Ethics Statement}

% N/A
% \section*{Acknowledgements}
We thank Joshua Howland, Alex Polozov, and Shrestha Basu Mallick for kindly sharing IPython notebook data. We thank the Google Pragma team for comments and discussions, Massimo Nicosia and William Cohen for their valuable feedback. We thank Ahmed Masry for providing us with ChartQA related data and baseline model outputs.



% Entries for the entire Anthology, followed by custom entries
\bibliography{anthology,custom}
\bibliographystyle{acl_natbib}

%\clearpage

\appendix



\section{More Details on Datasets Used}\label{sec:dataset_details}

\paragraph{Chart-table pairs from the web.} The data was originally collected by \citet{masry-etal-2022-chartqa} and came from the below four sources:
\begin{compactitem}
    \item Statista: \url{www.statista.com}
    \item Pew: \url{www.pewresearch.org}
    \item Our World in Data: \url{ourworldindata.org}
    \item OECD: \url{www.oecd.org}
\end{compactitem}

\paragraph{Modules of MATH questions included.} We exclude overly complex math questions and only select the basic modules that would help with numerical reasoning. They are from the two areas of Arithmetic and Comparison. The individual modules included are
\begin{compactitem}
    \item Arithmetic
    \begin{compactitem}
        \item \texttt{add\_or\_sub}
        \item \texttt{add\_sub\_multiple}
        \item \texttt{div}
        \item \texttt{mixed}
        \item \texttt{mul}
        \item \texttt{mul\_div\_multiple}
    \end{compactitem}
     
    \item Comparison
    \begin{compactitem}
        \item \texttt{closest}
        \item \texttt{closest\_composed}
        \item \texttt{kth\_biggest}
        \item \texttt{kth\_biggest\_composed}
        \item \texttt{pair}
        \item \texttt{pair\_composed}
        \item \texttt{sort}
        \item \texttt{sort\_composed}
    \end{compactitem}
\end{compactitem}

Please see \citet{saxton2019analysing} for detailed descriptions about each module and how they are generated.


\section{Details of Baselines}\label{sec:baselines}
We introduce below the details of the baselines used in \Cref{tab:main}.

T5 is an encode-decoder Transformer model proposed by \citet{raffel2020exploring}. The baseline model T5 takes the concatenation of a linearized table (and a query, when the task is QA) as input, and aims to decode the target (answer or summarization). When the gold table is availible, the gold table is used as the input and the chart image is not used directly. VL-T5 proposed by \citet{cho2021unifying} is similar to T5 but also takes a visual input (i.e., the chart image) on the encoder side. VisionTaPas \citep{masry-etal-2022-chartqa} is modified from TaPas \citep{herzig-etal-2020-tapas} to incorporate the visual modality by adding a ViT model \citep{dosovitskiy2021image} and cross-modal fusion layers.
T5-OCR, VL-T5-OCR, and VisionTaPas-OCR are the same model as T5, VL-T5, and VisionTaPas, respectively. However, they do not assume the existence of gold table but use an OCR-based system to extract the data table from the chart image.
The above mentioned models and their performance numbers are all extracted from \citet{masry-etal-2022-chartqa} and \citet{kantharaj-etal-2022-chart}. Please see the original paper for more details.
Classification - Regression Chart Transformer (CRCT) \citep{levy2022classification} is the best performing model on PlotQA according to the PlotQA benchmark on \url{paperswithcode.com}. It uses a detector that extracts all textual and visual elements of chart then processes these elements with a multimodal Transformer.
%PaLI...

% \section{Additional Experiments and Analyses}\label{sec:more_exp}





\end{document}
